\chapter{Experimentos y resultados}
	\label{sec:experimentos}



\section{Limitaciones y Desafíos}
\label{sec:limitaciones_experimentos}



\subsection{Rigidez Arquitectónica según el orden del sistema}


\subsection{Condicionamiento Numérico y Precisión de Punto Flotante}
El segundo desafío crítico es la susceptibilidad de los operadores de proyección a la brecha de precisión numérica. Los solucionadores iterativos, especialmente aquellos que cruzan las fronteras de conos semidefinidos positivos mediante descomposición espectral, son extremadamente sensibles a los errores de truncamiento y redondeo. 

En los entornos de aprendizaje automático, el estándar de la industria es utilizar precisión de punto flotante de 32 bits para maximizar la velocidad de cálculo en hardware paralelo. Sin embargo, bajo esta resolución, el algoritmo inducía inestabilidad numérica: matrices que teóricamente eran definidas positivas perdían dicha propiedad a nivel microscópico, generando valores indefinidos (NaNs) y colapsando el entrenamiento. Para mitigar este fenómeno, fue mandatorio forzar todo el flujo de datos computacionales —incluyendo los tensores de las plantas físicas extraídas y los pesos internos de la red neuronal— hacia una \textbf{precisión doble} (64 bits). Esta decisión sacrificó velocidad de procesamiento a cambio de garantizar la factibilidad de la matriz geométrica del sistema.

\subsection{Complejidad de Memoria y el Truncamiento del Desenrollado}
El tercer obstáculo, y posiblemente el más restrictivo, proviene del pase hacia atrás (\textit{backward pass}). El entrenamiento utiliza la técnica de desenrollado algorítmico de grafos computacionales para derivar a través del proceso de solución iterativa de la Desigualdad Matricial Lineal. Durante la fase de inferencia operativa, para asegurar que la matriz proyectada alcance la factibilidad absoluta, el algoritmo de Douglas-Rachford requiere iterar profundamente, alcanzando ejecuciones del orden de miles de ciclos (ej., 5000 iteraciones).

No obstante, almacenar el grafo computacional completo de 5000 operaciones secuenciales para calcular los gradientes provoca un agotamiento de la memoria de acceso aleatorio (RAM) del sistema gráfico. Para evitar este límite de hardware, el desenrollado de la proyección durante el entrenamiento debió truncarse a un número reducido de ciclos (típicamente 30 iteraciones).

\subsection{Sensibilidad del Margen $\epsilon$ e Inicialización Subóptima}
El truncamiento del pase hacia atrás genera una discrepancia algorítmica fundamental conocida como el "Gap de Entrenamiento". La red neuronal aprende a predecir basándose en el gradiente de solo 30 iteraciones, pero el sistema real se evalúa bajo una convergencia profunda. Esto provoca que la red entregue una \textbf{inicialización subóptima} para el solucionador en la etapa de proyección final.

Este problema se exacerba debido a la rigidez del hiperparámetro de seguridad $\epsilon$. Para sistemas dinámicos "marginales" —aquellos donde los polos de la planta $A_i$ se encuentran inherentemente al borde de la inestabilidad— la intersección geométrica entre el espacio afín y el cono semidefinido positivo se vuelve extremadamente delgada. Una inicialización subóptima combinada con un margen $\epsilon$ estricto hace que la convergencia iterativa sea altamente ineficiente, requiriendo un número masivo de iteraciones adicionales para que la matriz logre ingresar a la región matemáticamente segura. 